% --- Archon physics formalization source begin ---
% archon:physics
% archon:covers PhyXMiniProblems/problem_phyx_mini_0643.lean
% archon:source-report reports/phyx_mini/problem_phyx_mini_0643.source.json

\chapter{Physics problem phyx\_mini\_0643}
\label{ch:PhyXMiniProblems_problem_phyx_mini_0643}

\paragraph{Problem source.}
\#\# Physical scenario

An experiment was performed in which neutrons were shot through two slits spaced \$0.10\$ \$nm\$ apart and detected \$3.5\$ \$m\$ behind the slits. Figure shows the detector output. Notice the \$100\$ \$\textbackslash{}mu m\$ scale on the figure.

\#\# Question

To one significant figure, what was the speed of the neutrons?

\#\# Answer choices

- A: A: 150 m/s

- B: B: 200 m/s

- C: C: 250 m/s

- D: D: 300 m/s

\#\# Figure caption (auxiliary; use the image as primary evidence)

The image is a line graph illustrating neutron intensity. The horizontal axis is labeled with a scale showing a distance of "100 µm." The vertical axis represents "Neutron intensity." The graph depicts a wave-like pattern with multiple peaks and valleys, starting and ending at lower points with varying heights and intervals in between. The pattern indicates fluctuations in neutron intensity across a specific distance.

\#\# Dataset metadata

Quantum Phenomena; Physical Model Grounding Reasoning, Multi-Formula Reasoning

\paragraph{Recorded answer/context.}
B: B: 200 m/s

\paragraph{Figure/image path.}
/root/proposal\_for\_physic/science-mango/phyx\_mini\_run/phyx\_data/test\_image/643.png

\paragraph{Formalization target.}
create a compiling Lean file with sorry bodies at `PhyXMiniProblems/problem\_phyx\_mini\_0643.lean`. The Lean declarations must preserve the physical quantities, dimensions or dimensional roles, figure labels, governing-law hypotheses, and final relation expressed by this problem.
Use Mathlib/Physlib names found through LeanExplore where available. If a physics API is missing, introduce faithful local abstractions rather than scalar placeholder aliases.

\begin{theorem}[Physics formalization target]
\label{thm:physics:phyx_mini_0643:target}
\lean{PhyXMiniProblems.ProblemPhyXMini0643.problem_phyx_mini_0643}
\uses{def:physics:phyx-mini-0643:phyxminiproblems-problemphyxmini0643-lengthinfemtometers, def:physics:phyx-mini-0643:phyxminiproblems-problemphyxmini0643-lengthinmicrometers, def:physics:phyx-mini-0643:phyxminiproblems-problemphyxmini0643-speedinmeterspersecond, def:physics:phyx-mini-0643:phyxminiproblems-problemphyxmini0643-neutrondoubleslitsetup, def:physics:phyx-mini-0643:phyxminiproblems-problemphyxmini0643-matchesneutrondoubleslitscenario, def:physics:phyx-mini-0643:phyxminiproblems-problemphyxmini0643-matchesproblemlengthreadouts, def:physics:phyx-mini-0643:phyxminiproblems-problemphyxmini0643-matchessupplieddetectorfigure, def:physics:phyx-mini-0643:phyxminiproblems-problemphyxmini0643-usesstandardneutronreferencedata, def:physics:phyx-mini-0643:phyxminiproblems-problemphyxmini0643-hasphysicalneutroninterferenceparameters, def:physics:phyx-mini-0643:phyxminiproblems-problemphyxmini0643-satisfiesneutroninterferencelaws, def:physics:phyx-mini-0643:phyxminiproblems-problemphyxmini0643-answerchoice, def:physics:phyx-mini-0643:phyxminiproblems-problemphyxmini0643-roundstoonesignificantfigure, def:physics:phyx-mini-0643:phyxminiproblems-problemphyxmini0643-matchesdisplayedspeed}
The assigned autoformalize agent should translate this physics problem into Lean declarations in the covered file, with theorem and lemma proof bodies written as `by sorry`.
\end{theorem}
\begin{proof}
This is an autoformalization task, not a proof task. Produce statements that can later be proved by the physics prover without weakening the physical model.
\end{proof}

% --- Archon declaration topology begin ---
\section{Declaration topology}
\begin{definition}[Length Quantity]
  \label{def:physics:phyx-mini-0643:phyxminiproblems-problemphyxmini0643-lengthquantity}
  \lean{PhyXMiniProblems.ProblemPhyXMini0643.LengthQuantity}
  A nonnegative, unit-independent physical length
\end{definition}
\begin{proof}
  This is the declared model definition of Length Quantity.
\end{proof}
\begin{definition}[Mass Quantity]
  \label{def:physics:phyx-mini-0643:phyxminiproblems-problemphyxmini0643-massquantity}
  \lean{PhyXMiniProblems.ProblemPhyXMini0643.MassQuantity}
  A nonnegative, unit-independent physical mass
\end{definition}
\begin{proof}
  This is the declared model definition of Mass Quantity.
\end{proof}
\begin{definition}[action Dimension]
  \label{def:physics:phyx-mini-0643:phyxminiproblems-problemphyxmini0643-actiondimension}
  \lean{PhyXMiniProblems.ProblemPhyXMini0643.actionDimension}
  The dimension of action, M L² T⁻¹, carried by Planck's constant
\end{definition}
\begin{proof}
  This is the declared model definition of action Dimension.
\end{proof}
\begin{definition}[Action Quantity]
  \label{def:physics:phyx-mini-0643:phyxminiproblems-problemphyxmini0643-actionquantity}
  \lean{PhyXMiniProblems.ProblemPhyXMini0643.ActionQuantity}
  \uses{def:physics:phyx-mini-0643:phyxminiproblems-problemphyxmini0643-actiondimension}
  A nonnegative, unit-independent physical action
\end{definition}
\begin{proof}
  This is the declared model definition of Action Quantity.
\end{proof}
\begin{definition}[length Readout]
  \label{def:physics:phyx-mini-0643:phyxminiproblems-problemphyxmini0643-lengthreadout}
  \lean{PhyXMiniProblems.ProblemPhyXMini0643.lengthReadout}
  \uses{def:physics:phyx-mini-0643:phyxminiproblems-problemphyxmini0643-lengthquantity}
  Read a physical length in a selected named length unit
\end{definition}
\begin{proof}
  This is the declared model definition of length Readout.
\end{proof}
\begin{definition}[length In Meters]
  \label{def:physics:phyx-mini-0643:phyxminiproblems-problemphyxmini0643-lengthinmeters}
  \lean{PhyXMiniProblems.ProblemPhyXMini0643.lengthInMeters}
  \uses{def:physics:phyx-mini-0643:phyxminiproblems-problemphyxmini0643-lengthquantity, def:physics:phyx-mini-0643:phyxminiproblems-problemphyxmini0643-lengthreadout}
  Metre readout of a physical length
\end{definition}
\begin{proof}
  This is the declared model definition of length In Meters.
\end{proof}
\begin{definition}[length In Femtometers]
  \label{def:physics:phyx-mini-0643:phyxminiproblems-problemphyxmini0643-lengthinfemtometers}
  \lean{PhyXMiniProblems.ProblemPhyXMini0643.lengthInFemtometers}
  \uses{def:physics:phyx-mini-0643:phyxminiproblems-problemphyxmini0643-lengthquantity, def:physics:phyx-mini-0643:phyxminiproblems-problemphyxmini0643-lengthreadout}
  Femtometre readout of a physical length
\end{definition}
\begin{proof}
  This is the declared model definition of length In Femtometers.
\end{proof}
\begin{definition}[length In Nanometers]
  \label{def:physics:phyx-mini-0643:phyxminiproblems-problemphyxmini0643-lengthinnanometers}
  \lean{PhyXMiniProblems.ProblemPhyXMini0643.lengthInNanometers}
  \uses{def:physics:phyx-mini-0643:phyxminiproblems-problemphyxmini0643-lengthquantity, def:physics:phyx-mini-0643:phyxminiproblems-problemphyxmini0643-lengthreadout}
  Nanometre readout of a physical length
\end{definition}
\begin{proof}
  This is the declared model definition of length In Nanometers.
\end{proof}
\begin{definition}[length In Micrometers]
  \label{def:physics:phyx-mini-0643:phyxminiproblems-problemphyxmini0643-lengthinmicrometers}
  \lean{PhyXMiniProblems.ProblemPhyXMini0643.lengthInMicrometers}
  \uses{def:physics:phyx-mini-0643:phyxminiproblems-problemphyxmini0643-lengthquantity, def:physics:phyx-mini-0643:phyxminiproblems-problemphyxmini0643-lengthreadout}
  Micrometre readout of a physical length
\end{definition}
\begin{proof}
  This is the declared model definition of length In Micrometers.
\end{proof}
\begin{definition}[mass In Kilograms]
  \label{def:physics:phyx-mini-0643:phyxminiproblems-problemphyxmini0643-massinkilograms}
  \lean{PhyXMiniProblems.ProblemPhyXMini0643.massInKilograms}
  \uses{def:physics:phyx-mini-0643:phyxminiproblems-problemphyxmini0643-massquantity}
  Kilogram readout of a physical mass
\end{definition}
\begin{proof}
  This is the declared model definition of mass In Kilograms.
\end{proof}
\begin{definition}[speed In Meters Per Second]
  \label{def:physics:phyx-mini-0643:phyxminiproblems-problemphyxmini0643-speedinmeterspersecond}
  \lean{PhyXMiniProblems.ProblemPhyXMini0643.speedInMetersPerSecond}
  Metre-per-second readout of a physical speed
\end{definition}
\begin{proof}
  This is the declared model definition of speed In Meters Per Second.
\end{proof}
\begin{definition}[action In Joule Seconds]
  \label{def:physics:phyx-mini-0643:phyxminiproblems-problemphyxmini0643-actioninjouleseconds}
  \lean{PhyXMiniProblems.ProblemPhyXMini0643.actionInJouleSeconds}
  \uses{def:physics:phyx-mini-0643:phyxminiproblems-problemphyxmini0643-actionquantity}
  Coherent-SI readout of an action, in joule-seconds
\end{definition}
\begin{proof}
  This is the declared model definition of action In Joule Seconds.
\end{proof}
\begin{definition}[Particle Species]
  \label{def:physics:phyx-mini-0643:phyxminiproblems-problemphyxmini0643-particlespecies}
  \lean{PhyXMiniProblems.ProblemPhyXMini0643.ParticleSpecies}
  Particle species sent through the slits
\end{definition}
\begin{proof}
  This is the declared model definition of Particle Species.
\end{proof}
\begin{definition}[Slit Arrangement]
  \label{def:physics:phyx-mini-0643:phyxminiproblems-problemphyxmini0643-slitarrangement}
  \lean{PhyXMiniProblems.ProblemPhyXMini0643.SlitArrangement}
  Aperture arrangement used by the experiment
\end{definition}
\begin{proof}
  This is the declared model definition of Slit Arrangement.
\end{proof}
\begin{definition}[Propagation Regime]
  \label{def:physics:phyx-mini-0643:phyxminiproblems-problemphyxmini0643-propagationregime}
  \lean{PhyXMiniProblems.ProblemPhyXMini0643.PropagationRegime}
  Approximation regime used by the textbook calculation
\end{definition}
\begin{proof}
  This is the declared model definition of Propagation Regime.
\end{proof}
\begin{definition}[Plot Axis]
  \label{def:physics:phyx-mini-0643:phyxminiproblems-problemphyxmini0643-plotaxis}
  \lean{PhyXMiniProblems.ProblemPhyXMini0643.PlotAxis}
  The two axes of the detector-output graph
\end{definition}
\begin{proof}
  This is the declared model definition of Plot Axis.
\end{proof}
\begin{definition}[Plot Axis Role]
  \label{def:physics:phyx-mini-0643:phyxminiproblems-problemphyxmini0643-plotaxisrole}
  \lean{PhyXMiniProblems.ProblemPhyXMini0643.PlotAxisRole}
  Physical role assigned to an axis of the detector-output graph
\end{definition}
\begin{proof}
  This is the declared model definition of Plot Axis Role.
\end{proof}
\begin{definition}[Plot Label]
  \label{def:physics:phyx-mini-0643:phyxminiproblems-problemphyxmini0643-plotlabel}
  \lean{PhyXMiniProblems.ProblemPhyXMini0643.PlotLabel}
  Printed labels whose presence is relevant to the detector graph
\end{definition}
\begin{proof}
  This is the declared model definition of Plot Label.
\end{proof}
\begin{definition}[Figure Landmark]
  \label{def:physics:phyx-mini-0643:phyxminiproblems-problemphyxmini0643-figurelandmark}
  \lean{PhyXMiniProblems.ProblemPhyXMini0643.FigureLandmark}
  Horizontal landmarks measured from the supplied raster
\end{definition}
\begin{proof}
  This is the declared model definition of Figure Landmark.
\end{proof}
\begin{definition}[Neutron Detector Figure]
  \label{def:physics:phyx-mini-0643:phyxminiproblems-problemphyxmini0643-neutrondetectorfigure}
  \lean{PhyXMiniProblems.ProblemPhyXMini0643.NeutronDetectorFigure}
  \uses{def:physics:phyx-mini-0643:phyxminiproblems-problemphyxmini0643-lengthquantity, def:physics:phyx-mini-0643:phyxminiproblems-problemphyxmini0643-plotaxis, def:physics:phyx-mini-0643:phyxminiproblems-problemphyxmini0643-plotaxisrole, def:physics:phyx-mini-0643:phyxminiproblems-problemphyxmini0643-plotlabel, def:physics:phyx-mini-0643:phyxminiproblems-problemphyxmini0643-figurelandmark}
  Presentation-level information from the primary figure. Pixel coordinates are dimensionless raster measurements, whereas scaleBarLength is a genuine physical detector length
\end{definition}
\begin{proof}
  This is the declared model definition of Neutron Detector Figure.
\end{proof}
\begin{definition}[Neutron Double Slit Setup]
  \label{def:physics:phyx-mini-0643:phyxminiproblems-problemphyxmini0643-neutrondoubleslitsetup}
  \lean{PhyXMiniProblems.ProblemPhyXMini0643.NeutronDoubleSlitSetup}
  \uses{def:physics:phyx-mini-0643:phyxminiproblems-problemphyxmini0643-lengthquantity, def:physics:phyx-mini-0643:phyxminiproblems-problemphyxmini0643-massquantity, def:physics:phyx-mini-0643:phyxminiproblems-problemphyxmini0643-actionquantity, def:physics:phyx-mini-0643:phyxminiproblems-problemphyxmini0643-particlespecies, def:physics:phyx-mini-0643:phyxminiproblems-problemphyxmini0643-slitarrangement, def:physics:phyx-mini-0643:phyxminiproblems-problemphyxmini0643-propagationregime, def:physics:phyx-mini-0643:phyxminiproblems-problemphyxmini0643-neutrondetectorfigure}
  Independent physical quantities in the experiment. The inferred adjacent fringe spacing, matter wavelength, and neutron speed are fields rather than definitions; the calibration and governing-law premises below relate them
\end{definition}
\begin{proof}
  This is the declared model definition of Neutron Double Slit Setup.
\end{proof}
\begin{definition}[Matches Neutron Double Slit Scenario]
  \label{def:physics:phyx-mini-0643:phyxminiproblems-problemphyxmini0643-matchesneutrondoubleslitscenario}
  \lean{PhyXMiniProblems.ProblemPhyXMini0643.MatchesNeutronDoubleSlitScenario}
  \uses{def:physics:phyx-mini-0643:phyxminiproblems-problemphyxmini0643-neutrondoubleslitsetup}
  Qualitative physical scenario and approximation selected by the problem
\end{definition}
\begin{proof}
  This is the declared model definition of Matches Neutron Double Slit Scenario.
\end{proof}
\begin{definition}[Matches Problem Length Readouts]
  \label{def:physics:phyx-mini-0643:phyxminiproblems-problemphyxmini0643-matchesproblemlengthreadouts}
  \lean{PhyXMiniProblems.ProblemPhyXMini0643.MatchesProblemLengthReadouts}
  \uses{def:physics:phyx-mini-0643:phyxminiproblems-problemphyxmini0643-lengthinmeters, def:physics:phyx-mini-0643:phyxminiproblems-problemphyxmini0643-lengthinnanometers, def:physics:phyx-mini-0643:phyxminiproblems-problemphyxmini0643-neutrondoubleslitsetup}
  The two calibrated physical lengths stated in the prose. This premise contains no fringe spacing, wavelength, neutron speed, or answer choice
\end{definition}
\begin{proof}
  This is the declared model definition of Matches Problem Length Readouts.
\end{proof}
\begin{definition}[Matches Supplied Detector Figure]
  \label{def:physics:phyx-mini-0643:phyxminiproblems-problemphyxmini0643-matchessupplieddetectorfigure}
  \lean{PhyXMiniProblems.ProblemPhyXMini0643.MatchesSuppliedDetectorFigure}
  \uses{def:physics:phyx-mini-0643:phyxminiproblems-problemphyxmini0643-lengthinmicrometers, def:physics:phyx-mini-0643:phyxminiproblems-problemphyxmini0643-neutrondetectorfigure}
  Rounded landmark coordinates measured from image 643 and the physical label on its scale bar. The two central peak intervals are averaged to reduce the one-pixel asymmetry introduced by rasterization. These are figure/data readouts, not wavelength or speed conclusions
\end{definition}
\begin{proof}
  This is the declared model definition of Matches Supplied Detector Figure.
\end{proof}
\begin{definition}[Uses Standard Neutron Reference Data]
  \label{def:physics:phyx-mini-0643:phyxminiproblems-problemphyxmini0643-usesstandardneutronreferencedata}
  \lean{PhyXMiniProblems.ProblemPhyXMini0643.UsesStandardNeutronReferenceData}
  \uses{def:physics:phyx-mini-0643:phyxminiproblems-problemphyxmini0643-massinkilograms, def:physics:phyx-mini-0643:phyxminiproblems-problemphyxmini0643-actioninjouleseconds, def:physics:phyx-mini-0643:phyxminiproblems-problemphyxmini0643-neutrondoubleslitsetup}
  Standard reference data used by the calculation. Physlib provides the scalar reduced Planck constant Constants.ℏ in joule-seconds, so the ordinary Planck action has SI readout 2 π ℏ. Neither calibration contains the requested speed
\end{definition}
\begin{proof}
  This is the declared model definition of Uses Standard Neutron Reference Data.
\end{proof}
\begin{definition}[Has Physical Neutron Interference Parameters]
  \label{def:physics:phyx-mini-0643:phyxminiproblems-problemphyxmini0643-hasphysicalneutroninterferenceparameters}
  \lean{PhyXMiniProblems.ProblemPhyXMini0643.HasPhysicalNeutronInterferenceParameters}
  \uses{def:physics:phyx-mini-0643:phyxminiproblems-problemphyxmini0643-lengthinmeters, def:physics:phyx-mini-0643:phyxminiproblems-problemphyxmini0643-massinkilograms, def:physics:phyx-mini-0643:phyxminiproblems-problemphyxmini0643-speedinmeterspersecond, def:physics:phyx-mini-0643:phyxminiproblems-problemphyxmini0643-actioninjouleseconds, def:physics:phyx-mini-0643:phyxminiproblems-problemphyxmini0643-neutrondoubleslitsetup}
  Positivity conditions selecting a physically meaningful experiment
\end{definition}
\begin{proof}
  This is the declared model definition of Has Physical Neutron Interference Parameters.
\end{proof}
\begin{definition}[Satisfies Neutron Interference Laws]
  \label{def:physics:phyx-mini-0643:phyxminiproblems-problemphyxmini0643-satisfiesneutroninterferencelaws}
  \lean{PhyXMiniProblems.ProblemPhyXMini0643.SatisfiesNeutronInterferenceLaws}
  \uses{def:physics:phyx-mini-0643:phyxminiproblems-problemphyxmini0643-lengthreadout, def:physics:phyx-mini-0643:phyxminiproblems-problemphyxmini0643-lengthinmeters, def:physics:phyx-mini-0643:phyxminiproblems-problemphyxmini0643-massinkilograms, def:physics:phyx-mini-0643:phyxminiproblems-problemphyxmini0643-speedinmeterspersecond, def:physics:phyx-mini-0643:phyxminiproblems-problemphyxmini0643-actioninjouleseconds, def:physics:phyx-mini-0643:phyxminiproblems-problemphyxmini0643-neutrondoubleslitsetup}
  The three general governing relations used in the calculation: linear image calibration converts the average of the two central peak intervals into a physical detector length using the 100 μm bar; adjacent paraxial double-slit maxima obey Δy d = L λ; nonrelativistic neutrons obey de Broglie's relation λ m v = h. The first two relations are stated in arbitrary length units. They do not specialize any unknown to the result requested by this problem
\end{definition}
\begin{proof}
  This is the declared model definition of Satisfies Neutron Interference Laws.
\end{proof}
\begin{definition}[Answer Choice]
  \label{def:physics:phyx-mini-0643:phyxminiproblems-problemphyxmini0643-answerchoice}
  \lean{PhyXMiniProblems.ProblemPhyXMini0643.AnswerChoice}
  Labels of the four speed choices supplied with the problem
\end{definition}
\begin{proof}
  This is the declared model definition of Answer Choice.
\end{proof}
\begin{definition}[displayed Speed In Meters Per Second]
  \label{def:physics:phyx-mini-0643:phyxminiproblems-problemphyxmini0643-displayedspeedinmeterspersecond}
  \lean{PhyXMiniProblems.ProblemPhyXMini0643.displayedSpeedInMetersPerSecond}
  \uses{def:physics:phyx-mini-0643:phyxminiproblems-problemphyxmini0643-answerchoice}
  Speed in metres per second printed beside each answer label
\end{definition}
\begin{proof}
  This is the declared model definition of displayed Speed In Meters Per Second.
\end{proof}
\begin{definition}[recorded Dataset Answer]
  \label{def:physics:phyx-mini-0643:phyxminiproblems-problemphyxmini0643-recordeddatasetanswer}
  \lean{PhyXMiniProblems.ProblemPhyXMini0643.recordedDatasetAnswer}
  \uses{def:physics:phyx-mini-0643:phyxminiproblems-problemphyxmini0643-answerchoice}
  Answer label recorded by the source dataset, retained as metadata only
\end{definition}
\begin{proof}
  This is the declared model definition of recorded Dataset Answer.
\end{proof}
\begin{definition}[Rounds To One Significant Figure]
  \label{def:physics:phyx-mini-0643:phyxminiproblems-problemphyxmini0643-roundstoonesignificantfigure}
  \lean{PhyXMiniProblems.ProblemPhyXMini0643.RoundsToOneSignificantFigure}
  reported is value rounded to one significant figure when it consists of one nonzero decimal digit times a power of ten and ordinary nearest-place rounding of value gives that number
\end{definition}
\begin{proof}
  This is the declared model definition of Rounds To One Significant Figure.
\end{proof}
\begin{definition}[Matches Displayed Speed]
  \label{def:physics:phyx-mini-0643:phyxminiproblems-problemphyxmini0643-matchesdisplayedspeed}
  \lean{PhyXMiniProblems.ProblemPhyXMini0643.MatchesDisplayedSpeed}
  \uses{def:physics:phyx-mini-0643:phyxminiproblems-problemphyxmini0643-speedinmeterspersecond, def:physics:phyx-mini-0643:phyxminiproblems-problemphyxmini0643-neutrondoubleslitsetup, def:physics:phyx-mini-0643:phyxminiproblems-problemphyxmini0643-answerchoice, def:physics:phyx-mini-0643:phyxminiproblems-problemphyxmini0643-displayedspeedinmeterspersecond, def:physics:phyx-mini-0643:phyxminiproblems-problemphyxmini0643-roundstoonesignificantfigure}
  A displayed choice agrees with the modeled speed to one significant figure
\end{definition}
\begin{proof}
  This is the declared model definition of Matches Displayed Speed.
\end{proof}
\begin{lemma}[adjacent fringe spacing micrometers eq]
  \label{lem:physics:phyx-mini-0643:phyxminiproblems-problemphyxmini0643-adjacent-fringe-spacing-micrometers-eq}
  \lean{PhyXMiniProblems.ProblemPhyXMini0643.adjacent_fringe_spacing_micrometers_eq}
  \uses{def:physics:phyx-mini-0643:phyxminiproblems-problemphyxmini0643-lengthinmicrometers, def:physics:phyx-mini-0643:phyxminiproblems-problemphyxmini0643-neutrondoubleslitsetup, def:physics:phyx-mini-0643:phyxminiproblems-problemphyxmini0643-matchessupplieddetectorfigure, def:physics:phyx-mini-0643:phyxminiproblems-problemphyxmini0643-satisfiesneutroninterferencelaws}
  The primary-image calibration gives an averaged adjacent fringe spacing 100 μm * 111 / (2 * 84) = 925/14 μm
\end{lemma}
\begin{proof}
  The conclusion follows from the governing relations and figure readouts named in the statement of adjacent fringe spacing micrometers eq.
\end{proof}
\begin{lemma}[matter wavelength femtometers eq]
  \label{lem:physics:phyx-mini-0643:phyxminiproblems-problemphyxmini0643-matter-wavelength-femtometers-eq}
  \lean{PhyXMiniProblems.ProblemPhyXMini0643.matter_wavelength_femtometers_eq}
  \uses{def:physics:phyx-mini-0643:phyxminiproblems-problemphyxmini0643-lengthinfemtometers, def:physics:phyx-mini-0643:phyxminiproblems-problemphyxmini0643-neutrondoubleslitsetup, def:physics:phyx-mini-0643:phyxminiproblems-problemphyxmini0643-matchesproblemlengthreadouts, def:physics:phyx-mini-0643:phyxminiproblems-problemphyxmini0643-matchessupplieddetectorfigure, def:physics:phyx-mini-0643:phyxminiproblems-problemphyxmini0643-satisfiesneutroninterferencelaws}
  Combining that calibrated fringe spacing with d = 0.10 nm, L = 3.5 m, and the paraxial double-slit law gives λ = 185/98 fm
\end{lemma}
\begin{proof}
  The conclusion follows from the governing relations and figure readouts named in the statement of matter wavelength femtometers eq.
\end{proof}
\begin{lemma}[neutron speed meters per second eq]
  \label{lem:physics:phyx-mini-0643:phyxminiproblems-problemphyxmini0643-neutron-speed-meters-per-second-eq}
  \lean{PhyXMiniProblems.ProblemPhyXMini0643.neutron_speed_meters_per_second_eq}
  \uses{def:physics:phyx-mini-0643:phyxminiproblems-problemphyxmini0643-speedinmeterspersecond, def:physics:phyx-mini-0643:phyxminiproblems-problemphyxmini0643-neutrondoubleslitsetup, def:physics:phyx-mini-0643:phyxminiproblems-problemphyxmini0643-matchesproblemlengthreadouts, def:physics:phyx-mini-0643:phyxminiproblems-problemphyxmini0643-matchessupplieddetectorfigure, def:physics:phyx-mini-0643:phyxminiproblems-problemphyxmini0643-usesstandardneutronreferencedata, def:physics:phyx-mini-0643:phyxminiproblems-problemphyxmini0643-hasphysicalneutroninterferenceparameters, def:physics:phyx-mini-0643:phyxminiproblems-problemphyxmini0643-satisfiesneutroninterferencelaws}
  Substituting the image-derived wavelength, the standard neutron mass, and ordinary Planck action into the de Broglie law gives the exact SI expression below. In particular, the source's literal 0.10 nm datum makes the result about 2.1 × 10\textasciicircum{}8 m/s, rather than any of the displayed speeds
\end{lemma}
\begin{proof}
  The conclusion follows from the governing relations and figure readouts named in the statement of neutron speed meters per second eq.
\end{proof}
\begin{lemma}[neutron speed in physically supported interval]
  \label{lem:physics:phyx-mini-0643:phyxminiproblems-problemphyxmini0643-neutron-speed-in-physically-supported-interval}
  \lean{PhyXMiniProblems.ProblemPhyXMini0643.neutron_speed_in_physically_supported_interval}
  \uses{def:physics:phyx-mini-0643:phyxminiproblems-problemphyxmini0643-speedinmeterspersecond, def:physics:phyx-mini-0643:phyxminiproblems-problemphyxmini0643-neutrondoubleslitsetup, def:physics:phyx-mini-0643:phyxminiproblems-problemphyxmini0643-matchesproblemlengthreadouts, def:physics:phyx-mini-0643:phyxminiproblems-problemphyxmini0643-matchessupplieddetectorfigure, def:physics:phyx-mini-0643:phyxminiproblems-problemphyxmini0643-usesstandardneutronreferencedata, def:physics:phyx-mini-0643:phyxminiproblems-problemphyxmini0643-hasphysicalneutroninterferenceparameters, def:physics:phyx-mini-0643:phyxminiproblems-problemphyxmini0643-satisfiesneutroninterferencelaws}
  The exact expression lies between 2.09 × 10\textasciicircum{}8 and 2.10 × 10\textasciicircum{}8 metres per second. This is the physically supported speed interval under the stated data and the textbook nonrelativistic de Broglie model
\end{lemma}
\begin{proof}
  The conclusion follows from the governing relations and figure readouts named in the statement of neutron speed in physically supported interval.
\end{proof}
\begin{lemma}[neutron speed rounds to two hundred million]
  \label{lem:physics:phyx-mini-0643:phyxminiproblems-problemphyxmini0643-neutron-speed-rounds-to-two-hundred-million}
  \lean{PhyXMiniProblems.ProblemPhyXMini0643.neutron_speed_rounds_to_two_hundred_million}
  \uses{def:physics:phyx-mini-0643:phyxminiproblems-problemphyxmini0643-speedinmeterspersecond, def:physics:phyx-mini-0643:phyxminiproblems-problemphyxmini0643-neutrondoubleslitsetup, def:physics:phyx-mini-0643:phyxminiproblems-problemphyxmini0643-matchesproblemlengthreadouts, def:physics:phyx-mini-0643:phyxminiproblems-problemphyxmini0643-matchessupplieddetectorfigure, def:physics:phyx-mini-0643:phyxminiproblems-problemphyxmini0643-usesstandardneutronreferencedata, def:physics:phyx-mini-0643:phyxminiproblems-problemphyxmini0643-hasphysicalneutroninterferenceparameters, def:physics:phyx-mini-0643:phyxminiproblems-problemphyxmini0643-satisfiesneutroninterferencelaws, def:physics:phyx-mini-0643:phyxminiproblems-problemphyxmini0643-roundstoonesignificantfigure}
  The modeled speed rounds to 2 × 10\textasciicircum{}8 m/s at one significant figure
\end{lemma}
\begin{proof}
  The conclusion follows from the governing relations and figure readouts named in the statement of neutron speed rounds to two hundred million.
\end{proof}
\begin{lemma}[no displayed speed matches]
  \label{lem:physics:phyx-mini-0643:phyxminiproblems-problemphyxmini0643-no-displayed-speed-matches}
  \lean{PhyXMiniProblems.ProblemPhyXMini0643.no_displayed_speed_matches}
  \uses{def:physics:phyx-mini-0643:phyxminiproblems-problemphyxmini0643-neutrondoubleslitsetup, def:physics:phyx-mini-0643:phyxminiproblems-problemphyxmini0643-matchesproblemlengthreadouts, def:physics:phyx-mini-0643:phyxminiproblems-problemphyxmini0643-matchessupplieddetectorfigure, def:physics:phyx-mini-0643:phyxminiproblems-problemphyxmini0643-usesstandardneutronreferencedata, def:physics:phyx-mini-0643:phyxminiproblems-problemphyxmini0643-hasphysicalneutroninterferenceparameters, def:physics:phyx-mini-0643:phyxminiproblems-problemphyxmini0643-satisfiesneutroninterferencelaws, def:physics:phyx-mini-0643:phyxminiproblems-problemphyxmini0643-answerchoice, def:physics:phyx-mini-0643:phyxminiproblems-problemphyxmini0643-matchesdisplayedspeed}
  Consequently none of the four displayed values 150, 200, 250, or 300 metres per second matches the speed inferred from the literal source data. The recorded label B remains available only through recordedDatasetAnswer
\end{lemma}
\begin{proof}
  The conclusion follows from the governing relations and figure readouts named in the statement of no displayed speed matches.
\end{proof}
% --- Archon declaration topology end ---
% --- Archon physics formalization source end ---
